% GENERATED FILE. Edit canonical lessons, then run npm run generate:books.
% canonical-source-hash: fnv1a64:e7a3a4b85f3f2bf1
% canonical-lessons: UR-C420-amdani, UR-C420-bajat, UR-C420-kifayat, UR-C420-bantna, UR-C420-bhunana

\chapter{Income and Budgets}
\label{ch:ur-income-and-budgets}
\hlchaptermodality{420}

\begin{chapteropening}
\textbf{By the end of this chapter:} \emph{I can talk about what comes in, what goes out and how to make it last.}

Complete the last lesson of chapter 420: I can talk about what comes in, what goes out and how to make it last.
\end{chapteropening}

\section[\texorpdfstring{\ur{آمدنی} (āmdanī)}{آمدنی (āmdanī)}]{\ur{آمدنی} (āmdanī) — an income}

\label{lesson:UR-C420-amdani}



\begin{quote}
Before the new one: say the Urdu for an asset, then the Urdu for to pack up.
\end{quote}

\subsection*{You'll want to know: \ur{آمدنی}}
\textbf{\ur{آمدنی}} — \emph{āmdanī} — \textquotedblleft{}an income\textquotedblright{}.

\begin{quote}
\ur{اس} \ur{کی} \ur{آمدنی} \ur{کم} \ur{ہے۔}
\end{quote}

\begin{quote}
\emph{us kī āmdanī kam hai.}
\end{quote}

\begin{quote}
He doesn't earn much.
\end{quote}

Feminine.

\begin{cousinweb}
\textbf{\ur{آمدنی}} (\emph{āmdanī}) is usually traced to Persian \emph{āmadanī}, \textquotedblleft{}what comes in\textquotedblright{}, from the verb for \textquotedblleft{}to come\textquotedblright{}.
\end{cousinweb}

\begin{grammarlens}[title={What you've built}]
One more word for getting around.
\end{grammarlens}

\subsection*{Your turn}
\begin{itemize}
\raggedright
  \item \emph{Say it:} \emph{āmdanī}
  \item \emph{Say it:} \emph{āmdanī}, once more
  \item \emph{Say it:} \emph{us kī āmdanī kam hai.} — He doesn't earn much.
  \item \emph{Recall:} say the Urdu for an asset, then the Urdu for to pack up, then say \emph{āmdanī} again
\end{itemize}

\begin{tcolorbox}[breakable,colback=teal!4,colframe=teal!35!black,title={Before you move on}]
What does \ur{آمدنی} mean? (\textquotedblleft{}An income\textquotedblright{}.) Say it once more.
\end{tcolorbox}
\section[\texorpdfstring{\ur{بجٹ} (bajaṭ)}{بجٹ (bajaṭ)}]{\ur{بجٹ} (bajaṭ) — a budget}

\label{lesson:UR-C420-bajat}



\begin{quote}
Before the new one: say the Urdu for to pack up, then the Urdu for an income.
\end{quote}

\subsection*{You'll want to know: \ur{بجٹ}}
\textbf{\ur{بجٹ}} — \emph{bajaṭ} — \textquotedblleft{}a budget\textquotedblright{}.

\begin{quote}
\ur{ہمارا} \ur{بجٹ} \ur{ہماری} \ur{آمدنی} \ur{پر} \ur{منحصر} \ur{ہے۔}
\end{quote}

\begin{quote}
\emph{hamārā bajaṭ hamārī āmdanī par munhasir hai.}
\end{quote}

\begin{quote}
Our budget depends on our income.
\end{quote}

Masculine.

\begin{cousinweb}
\textbf{\ur{بجٹ}} (\emph{bajaṭ}) is borrowed from English \emph{budget}.
\end{cousinweb}

\begin{grammarlens}[title={What you've built}]
One more word for getting around.
\end{grammarlens}

\subsection*{Your turn}
\begin{itemize}
\raggedright
  \item \emph{Say it:} \emph{bajaṭ}
  \item \emph{Say it:} \emph{bajaṭ}, once more
  \item \emph{Say it:} \emph{hamārā bajaṭ hamārī āmdanī par munhasir hai.} — Our budget depends on our income.
  \item \emph{Recall:} say the Urdu for to pack up, then the Urdu for an income, then say \emph{bajaṭ} again
\end{itemize}

\begin{tcolorbox}[breakable,colback=teal!4,colframe=teal!35!black,title={Before you move on}]
What does \ur{بجٹ} mean? (\textquotedblleft{}A budget\textquotedblright{}.) Say it once more.
\end{tcolorbox}
\section[\texorpdfstring{\ur{کفایت} (kifāyat)}{کفایت (kifāyat)}]{\ur{کفایت} (kifāyat) — thrift; economy}

\label{lesson:UR-C420-kifayat}



\begin{quote}
Before the new one: say the Urdu for an income, then the Urdu for a budget.
\end{quote}

\subsection*{You'll want to know: \ur{کفایت}}
\textbf{\ur{کفایت}} — \emph{kifāyat} — \textquotedblleft{}thrift; economy\textquotedblright{}.

\begin{quote}
\ur{ہمیں} \ur{کفایت} \ur{سے} \ur{خرچ} \ur{کرنا} \ur{چاہیے۔}
\end{quote}

\begin{quote}
\emph{hameṅ kifāyat se kharch karnā chāhiye.}
\end{quote}

\begin{quote}
We should spend carefully.
\end{quote}

Feminine.

\begin{cousinweb}
\textbf{\ur{کفایت}} (\emph{kifāyat}) is usually traced, through Persian, to Arabic \emph{kifāya}, \textquotedblleft{}sufficiency, enough\textquotedblright{}; \textbf{\ur{کفایت} \ur{سے}} is \textquotedblleft{}thriftily\textquotedblright{}.
\end{cousinweb}

\begin{grammarlens}[title={What you've built}]
One more word for getting around.
\end{grammarlens}

\subsection*{Your turn}
\begin{itemize}
\raggedright
  \item \emph{Say it:} \emph{kifāyat}
  \item \emph{Say it:} \emph{kifāyat}, once more
  \item \emph{Say it:} \emph{hameṅ kifāyat se kharch karnā chāhiye.} — We should spend carefully.
  \item \emph{Recall:} say the Urdu for an income, then the Urdu for a budget, then say \emph{kifāyat} again
\end{itemize}

\begin{tcolorbox}[breakable,colback=teal!4,colframe=teal!35!black,title={Before you move on}]
What does \ur{کفایت} mean? (\textquotedblleft{}Thrift; economy\textquotedblright{}.) Say it once more.
\end{tcolorbox}
\section[\texorpdfstring{\ur{بانٹنا} (bāṅṭnā)}{بانٹنا (bāṅṭnā)}]{\ur{بانٹنا} (bāṅṭnā) — to share out; to divide}

\label{lesson:UR-C420-bantna}



\begin{quote}
Before the new one: say the Urdu for a budget, then the Urdu for thrift.
\end{quote}

\subsection*{You'll want to know: \ur{بانٹنا}}
\textbf{\ur{بانٹنا}} — \emph{bāṅṭnā} — \textquotedblleft{}to share out; to divide\textquotedblright{}.

\begin{quote}
\ur{دونوں} \ur{بھائیوں} \ur{نے} \ur{منافع} \ur{بانٹا۔}
\end{quote}

\begin{quote}
\emph{donoṅ bhāiyoṅ ne munāfā bāṅṭā.}
\end{quote}

\begin{quote}
The two brothers shared the profit.
\end{quote}

Past, with \ur{نے} (he or she …): \ur{اس} \ur{نے} \ur{مٹھائی} \ur{بانٹی}, \emph{us ne miṭhāī bāṅṭī}.

\begin{cousinweb}
\textbf{\ur{بانٹنا}} (\emph{bāṅṭnā}) is inherited: CDIAL lists the Hindi cousin \emph{bā̃ṭnā}, \textquotedblleft{}shares\textquotedblright{}, under Sanskrit \emph{váṇṭati}, \textquotedblleft{}shares\textquotedblright{} (CDIAL 11238).
\end{cousinweb}

\begin{grammarlens}[title={What you've built}]
One more word for getting around.
\end{grammarlens}

\subsection*{Your turn}
\begin{itemize}
\raggedright
  \item \emph{Say it:} \emph{bāṅṭnā}
  \item \emph{Say it:} \emph{bāṅṭnā}, once more
  \item \emph{Say it:} \emph{donoṅ bhāiyoṅ ne munāfā bāṅṭā.} — The two brothers shared the profit.
  \item \emph{Recall:} say the Urdu for a budget, then the Urdu for thrift, then say \emph{bāṅṭnā} again
\end{itemize}

\begin{tcolorbox}[breakable,colback=teal!4,colframe=teal!35!black,title={Before you move on}]
What does \ur{بانٹنا} mean? (\textquotedblleft{}To share out; to divide\textquotedblright{}.) Say it once more.
\end{tcolorbox}
\section[\texorpdfstring{\ur{بھنانا} (bhunānā)}{بھنانا (bhunānā)}]{\ur{بھنانا} (bhunānā) — to cash (a cheque); to change (a note)}

\label{lesson:UR-C420-bhunana}



\begin{quote}
Before the new one: say the Urdu for thrift, then the Urdu for to share out.
\end{quote}

\subsection*{You'll want to know: \ur{بھنانا}}
\textbf{\ur{بھنانا}} — \emph{bhunānā} — \textquotedblleft{}to cash (a cheque); to change (a note)\textquotedblright{}.

\begin{quote}
\ur{سارا} \ur{نے} \ur{بینک} \ur{میں} \ur{چیک} \ur{بھنایا۔}
\end{quote}

\begin{quote}
\emph{sārā ne baink meṅ chek bhunāyā.}
\end{quote}

\begin{quote}
Sara cashed a cheque at the bank.
\end{quote}

Past, with \ur{نے} (he or she …): \ur{اس} \ur{نے} \ur{نوٹ} \ur{بھنایا}, \emph{us ne noṭ bhunāyā}.

\begin{cousinweb}
\textbf{\ur{بھنانا}} (\emph{bhunānā}) is inherited: CDIAL lists Hindi \emph{bhunnā}, \textquotedblleft{}to be changed (of money)\textquotedblright{}, and \emph{bhunānā} under Sanskrit \emph{bhajyate}, \textquotedblleft{}is broken\textquotedblright{} (CDIAL 9361): a big note is broken into small ones.
\end{cousinweb}

\begin{grammarlens}[title={What you've built}]
One more word for getting around.
\end{grammarlens}

\subsection*{Your turn}
\begin{itemize}
\raggedright
  \item \emph{Say it:} \emph{bhunānā}
  \item \emph{Say it:} \emph{bhunānā}, once more
  \item \emph{Say it:} \emph{sārā ne baink meṅ chek bhunāyā.} — Sara cashed a cheque at the bank.
  \item \emph{Recall:} say the Urdu for thrift, then the Urdu for to share out, then say \emph{bhunānā} again
\end{itemize}

\begin{tcolorbox}[breakable,colback=teal!4,colframe=teal!35!black,title={Before you move on}]
What does \ur{بھنانا} mean? (\textquotedblleft{}To cash (a cheque); to change (a note)\textquotedblright{}.) Say it once more.
\end{tcolorbox}
